\documentclass[11pt,a4paper]{article}

\usepackage{fontspec}
\setmainfont{Latin Modern Roman}
\setsansfont{Latin Modern Sans}
\setmonofont{Latin Modern Mono}
\usepackage[english]{babel}
\usepackage{amsmath,amssymb,amsthm}
\usepackage{algorithm}
\usepackage[noend]{algpseudocode}
\usepackage{booktabs}
\usepackage{multirow}
\usepackage{graphicx}
\usepackage[margin=2.5cm]{geometry}
\usepackage{xcolor}
\usepackage{enumitem}
\usepackage{caption}
\usepackage{subcaption}
\usepackage[hidelinks]{hyperref}
\hypersetup{
  pdftitle    = {Fast Greedy Heuristics with Linear-Time Redundant Removal
                 and Local Swap for the Minimum k-Dominating Set Problem
                 on Large-Scale Networks},
  pdfauthor   = {Tho Chi Luong},
  pdfsubject  = {Minimum k-dominating set problem},
  pdfkeywords = {k-dominating set; greedy heuristic; redundant removal;
                 local search; coverage count; large-scale networks},
  pdfcreator  = {},
  pdfproducer = {}
}

\newtheorem{theorem}{Theorem}
\newtheorem{lemma}[theorem]{Lemma}
\newtheorem{definition}[theorem]{Definition}
\newtheorem{proposition}[theorem]{Proposition}

\newcommand{\Nk}{\overline{N}_k}
\newcommand{\cov}{\mathrm{cov}}
\newcommand{\loss}{\mathrm{loss}}
\newcommand{\OPT}{\mathrm{OPT}}

% ── Title ────────────────────────────────────────────────────
\title{%
  Fast Greedy Heuristics with Linear-Time Redundant Removal\\
  and Local Swap for the Minimum $k$-Dominating Set Problem\\
  on Large-Scale Networks%
}

\author{
  Tho Chi Luong\thanks{Email: \texttt{tholc@vnu.edu.vn}}
  \\[6pt]
  International School,\\
  Vietnam National University, Hanoi\\
  Hanoi, Vietnam
}

\date{\today}

% =============================================================
\begin{document}
\maketitle

% ─────────────────────────────────────────────────────────────
% ABSTRACT
% ─────────────────────────────────────────────────────────────
\begin{abstract}
Given a graph $G=(V,E)$ and integer $k\ge 1$, a \emph{$k$-dominating
set} is a vertex subset $D\subseteq V$ inside whose
distance-$k$ neighbourhood every vertex of $G$ lies. The
minimum-cardinality version (M$k$DSP) is NP-hard and surfaces in
backbone selection for sensor networks, seed targeting for multi-hop
information cascades, and radius-bounded service placement.

This paper develops two post-optimisation primitives built around a
single shared data structure—the per-vertex \emph{coverage count}
$\cov(u)=|\{v\in D:u\in N_k(v)\}|$—and pairs them with the
multi-threshold construction GreedyTheta of
Nguyen et~al.~\cite{nguyen2020} to form heuristic pipelines on
large graphs. \emph{Fast Redundant Removal} (RR) exploits a single
invariant: $v\in D$ may be discarded whenever every vertex of
$N_k(v)$ already has coverage at least two. Maintaining
$\cov(\cdot)$ incrementally drops the cost of one sweep over $D$
from $O(|D|^2\,\Nk^2)$ to $O(|D|\,\Nk)$. \emph{Local Swap} (LS)
generalises RR by tolerating a small \emph{loss set} of
uniquely-covered vertices and replacing the offending $v$ with a
single substitute $w\notin D$ whose distance-$k$ ball contains that
loss. Both procedures preserve the invariant $\cov(u)\ge 1$ and
therefore output a valid $k$-dominating set at every step.

On nine fresh Erd\H{o}s--R\'enyi-like graphs with $n=100{,}000$,
$|E|\approx 2.5\times 10^5$ and average degree close to five (taking
$k=3$), the two-phase pipeline GT$+$RR trims the construction's
$|D|$ by \textbf{10.6\%} on average ($3{,}686\to 3{,}295$) for a mean
overhead of $1.4$~s (worst-case $4.0$\,s), and the three-phase
GT$+$LS$+$RR pipeline reaches a reduction of \textbf{30.1\%}
($3{,}686\to 2{,}577$) in roughly $185$~s per instance. RR is
essentially free; the local-search stage drives almost all of the
remaining gain.

\medskip\noindent
\textbf{Keywords:}
$k$-dominating set, multi-threshold greedy construction, linear-time
redundant removal, local search, coverage-count bookkeeping,
large-scale graphs.
\end{abstract}

% ─────────────────────────────────────────────────────────────
% VIETNAMESE ABSTRACT (yêu cầu của tạp chí trong nước)
% ─────────────────────────────────────────────────────────────
\begin{center}
\textbf{Heuristic tham lam với loại đỉnh dư thời gian tuyến tính
và hoán đổi cục bộ cho bài toán $k$-tập trội tối thiểu trên đồ thị
lớn}
\end{center}

\noindent\textbf{Tóm tắt.}
Cho đồ thị vô hướng $G=(V,E)$ và số nguyên $k\ge 1$, một $k$-tập
trội là tập đỉnh $D\subseteq V$ sao cho mọi đỉnh của $G$ đều cách
một đỉnh nào đó trong $D$ không quá $k$ cạnh. Bài toán
$k$-tập trội tối thiểu (M$k$DSP) thuộc lớp NP-khó và có mặt trong
nhiều ứng dụng: chọn xương sống cho mạng cảm biến không dây, gieo
hạt cho lan truyền thông tin trong mạng xã hội, và bố trí dịch vụ
với bán kính giới hạn. Bài báo đề xuất hai kỹ thuật hậu tối ưu
khai thác chung một bộ đếm bao phủ $\cov(u)=|\{v\in D:u\in N_k(v)\}|$:
(i) \emph{Loại đỉnh dư nhanh} (RR) chạy trong thời gian
$O(|D|\,\Nk)$, thay cho thuật toán cổ điển $O(|D|^2\,\Nk^2)$;
(ii) \emph{Hoán đổi cục bộ} (LS) mở rộng RR bằng phép $(1,1)$-swap
dựa trên tập mất bao phủ
$\loss(v)=\{u\in N_k(v):\cov(u)\le 1\}$, đổi $v$ lấy một đỉnh thay
thế $w\notin D$ với điều kiện $\loss(v)\subseteq N_k(w)$. Cả hai
thủ tục đều bảo toàn bất biến $\cov(u)\ge 1$, do đó luôn trả về
một $k$-tập trội hợp lệ.
Thực nghiệm trên 9~đồ thị ngẫu nhiên kiểu Erd\H{o}s--R\'enyi với
$n=100{,}000$ đỉnh, $|E|\approx 2{,}5\cdot 10^5$ và bậc trung bình
xấp xỉ năm (tham số $k=3$) cho thấy: đường ống GT$+$RR rút gọn
$|D|$ trung bình \textbf{10{,}6\%} ($3686\to 3295$) với phần phụ
trội trung bình $1{,}4$~s (lớn nhất $4{,}0$\,s); đường ống ba pha
GT$+$LS$+$RR đạt mức rút gọn \textbf{30{,}1\%}
($3686\to 2577$) trong khoảng $185$~giây cho mỗi đồ thị. RR gần
như không tốn chi phí; pha LS đóng góp hầu hết phần cải thiện.

\medskip\noindent
\textbf{Từ khóa.} bài toán $k$-tập trội; heuristic tham lam đa
ngưỡng; loại đỉnh dư thời gian tuyến tính; tìm kiếm cục bộ; bộ đếm
bao phủ; đồ thị lớn.

% =============================================================
\section{Introduction}
\label{sec:intro}
% =============================================================

\subsection{Motivation}

Domination problems ask one to anchor a graph by a small set of
vertices whose neighbourhoods cover everyone else. The generalisation
considered here, M$k$DSP, widens ``neighbourhood'' to mean the
distance-$k$ ball $N_k(\cdot)$ around a vertex; the case $k=1$ recovers
the classical minimum dominating set, whose hardness was established
already in~\cite{garey1979} and whose
approximability is bounded below by the
$\Omega(\log n)$ threshold of the set-cover reduction
through~\cite{raz1997}. The version with $k\ge 2$ inherits this
hardness while becoming more useful in practice: a single dominator
can now ``represent'' a much larger ball, so the resulting set is
typically several times smaller than the $k=1$ case.

Three settings motivate our study. The first is virtual-backbone
construction in wireless ad-hoc and sensor
networks~\cite{amis2000}, where increasing $k$ relaxes the cluster
radius and shrinks the number of cluster heads to be deployed. The
second is multi-hop seed selection for cascade processes on social
networks: classical influence-maximisation models, surveyed
in~\cite{kempe2003}, become an M$k$DSP instance once the threshold of
``successful infection'' is taken to be reachability within $k$
edges. The third is service deployment with a radius
budget~\cite{hassin1991}: a $k$-dominating set, on a road network
weighted by hop count, is precisely a feasible placement of facilities
whose distance to any user does not exceed $k$.

The technical obstacle in all three settings is the same: a
distance-$k$ search around every candidate vertex is expensive on
graphs with $10^5$ or more vertices, and exact ILP/MILP formulations
collapse well before that size. The state of the art for
million-vertex instances is the $\mathit{HEU_4}$ pipeline of
Nguyen et~al.~\cite{nguyen2020}: a $k$-isolated-cluster
preprocessor; a multi-threshold greedy phase identical to our GT; and
a partition-based MILP post-optimisation that calls CPLEX on each
partition. Table~4 of~\cite{nguyen2020} reports
$\mathit{HEU_4}$ wall-clock times between $1{,}695$ and
$42{,}600$~seconds for $k=3$ on real social networks of $89$K to $4$M
vertices, the MILP stage being the dominant cost and the only step
that requires a commercial solver licence.

Our question is whether one can match or exceed that solution
quality without any MILP, on a single thread, using only the
coverage-count primitive defined below.

\subsection{Contributions}

This paper contributes the following.

\begin{enumerate}[nosep]
\item \textbf{Linear-time redundant removal.}
  We isolate the coverage counter
  $\cov(u)\,=\,|\{v\in D:u\in N_k(v)\}|$
  as the single piece of state that makes redundancy detection cheap:
  a vertex $v\in D$ is redundant exactly when
  $\min_{u\in N_k(v)}\cov(u)\ge 2$. Computing $\cov(\cdot)$ once and
  maintaining it incrementally lets a full sweep over $D$ finish in
  $O(|D|\cdot\Nk)$ time, in place of the textbook
  $O(|D|^2\cdot\Nk^2)$.

\item \textbf{A $(1,1)$-swap built on the same primitive.}
  We push the coverage counter past the redundancy threshold: when
  $\loss(v)=\{u\in N_k(v):\cov(u)\le 1\}$ is non-empty but
  small ($|\loss(v)|\le L_{\max}$), LS hunts among at most $C_{\max}$
  outside candidates $w$ for one whose $N_k(w)$ covers $\loss(v)$.
  With $L_{\max}$ and $C_{\max}$ treated as constants the running time
  per pass is the same $O(|D|\cdot\Nk)$ as RR; we cap the number of
  passes at $P$ and obtain $O(P\cdot|D|\cdot\Nk)$ overall.

\item \textbf{One invariant, two correctness arguments.}
  We establish that both procedures preserve
  $\cov(u)\ge 1$ for all $u\in V$
  (Theorem~\ref{thm:rr-correct}, Lemma~\ref{lem:ls-correct}). Every
  intermediate set, and the final output, is therefore a valid
  $k$-dominating set.

\item \textbf{Two heuristic pipelines without MILP.}
  GT$+$RR (two phases) and GT$+$LS$+$RR (three phases) reuse the
  GreedyTheta scheduler of~\cite{nguyen2020} on the construction
  side but discard the MILP-based post-optimisation step of
  $\mathit{HEU_4}$ entirely. The $k$-isolated-cluster preprocessor
  is also dropped, since for $k=3$ it offers no measurable benefit on
  our benchmark.

\item \textbf{An empirical assessment.}
  On nine fresh random graphs with $n=100{,}000$,
  $|E|\!\approx\!2.5\!\times\!10^5$ and average degree close to five,
  taking $k=3$: GT$+$RR shrinks the GreedyTheta solution by
  $10.6\%$ on average; GT$+$LS$+$RR shrinks it by $30.1\%$. Adding RR
  to GT costs $1.4$~s on average (worst case $4.0$~s); the LS stage
  costs roughly $155$~s and explains essentially all of the additional
  gain. Detailed results are tabulated in Section~\ref{sec:experiments}.
\end{enumerate}

\subsection{Paper Organization}

Section~\ref{sec:prelim} fixes notation and the coverage-count and
loss-set primitives. Section~\ref{sec:related} reviews related work.
Section~\ref{sec:algo} presents GT, RR, LS, and the two pipelines, with
correctness and complexity proofs. Section~\ref{sec:experiments} reports
experimental results. Section~\ref{sec:discussion} discusses insights and
limitations. Section~\ref{sec:conclusion} concludes.


% =============================================================
\section{Preliminaries}
\label{sec:prelim}
% =============================================================

\subsection{Notation}

Throughout, $G=(V,E)$ is undirected, unweighted, and simple, with
$n=|V|$ and $m=|E|$. The unparameterised neighbourhood
$N(v)=\{u:(v,u)\in E\}$ and its closed version $N[v]=N(v)\cup\{v\}$
take their usual meaning; $\mathrm{dist}(u,v)$ is the number of edges
on a shortest $u$--$v$ path, with the convention $\mathrm{dist}(v,v)=0$.
The central object of this paper is the $k$-hop neighbourhood
\[
  N_k(v) \;=\; \{\,u\in V:\mathrm{dist}(u,v)\le k\,\},
\]
which includes $v$ itself. We write $\Nk$ for the average
$|N_k(\cdot)|$ over a context-dependent set: vertex-wise averages over
$V$ during the analysis of construction, and over $D$ during the
analysis of removal.

\subsection{Problem Definition}

\begin{definition}[M$k$DSP]
\label{def:mkdsp}
An instance of the \emph{minimum $k$-dominating set problem} is a pair
$(G,k)$. A subset $D\subseteq V$ is feasible if its distance-$k$ balls
jointly cover $V$, i.e.\ $\bigcup_{u\in D}N_k(u)=V$; equivalently, every
$v\in V$ admits a witness $u\in D$ with $\mathrm{dist}(u,v)\le k$.
The objective is to minimise $|D|$ over feasible subsets.
\end{definition}

For $k=1$, Definition~\ref{def:mkdsp} reduces to the classical minimum
dominating set, which is NP-complete in its decision version. The
reduction is preserved by the inclusion $\{1\}\subseteq\{k:k\ge 1\}$,
so M$k$DSP is NP-hard for every fixed $k\ge 1$.

\subsection{Coverage Count}

\begin{definition}[Coverage Count]
\label{def:cov}
For a candidate dominating set $D\subseteq V$, the \emph{coverage count}
of vertex $u\in V$ is
\[
  \cov(u,D) = |\{v\in D : u\in N_k(v)\}|
            = |\{v\in D : \mathrm{dist}(u,v)\le k\}|.
\]
\end{definition}

The set $D$ is a valid $k$-dominating set iff $\cov(u,D)\ge 1$ for every
$u\in V$. A vertex $v\in D$ is \emph{redundant} if $D\setminus\{v\}$ is
still valid; equivalently, $\cov(u,D)\ge 2$ for every $u\in N_k(v)$.

\subsection{Loss Set}

\begin{definition}[Loss Set]
\label{def:loss}
For a candidate $v\in D$, the \emph{loss set} is
\[
  \loss(v,D) = \{u\in N_k(v) : \cov(u,D) \le 1\}.
\]
\end{definition}

The loss set captures exactly those vertices that would become uncovered
were $v$ removed. Two extreme regimes guide the design of LS in
Section~\ref{sec:algo}:
\begin{itemize}[nosep]
  \item $\loss(v,D)=\emptyset$ — $v$ is redundant; remove unconditionally.
  \item $|\loss(v,D)|$ small — $v$ is \emph{near-redundant}; replace by a
    substitute $w\notin D$ with $\loss(v,D)\subseteq N_k(w)$.
  \item $|\loss(v,D)|$ large — keep $v$.
\end{itemize}


% =============================================================
\section{Related Work}
\label{sec:related}
% =============================================================

The literature surrounding M$k$DSP spans six threads: exact solvers,
greedy schedulers, post-optimisation, local search and metaheuristics,
sketch-based scalability, and learning-based methods. We treat each in
turn, then close with two short paragraphs on preprocessing and on
nearby variants. For broader background on graph domination as a
research area we point to the book by Haynes, Hedetniemi and
Slater~\cite{haynes1998}.

\subsection{Exact Methods}

The natural binary program for M$k$DSP introduces an indicator
$x_v\in\{0,1\}$ for membership in $D$ and one covering constraint per
vertex:
\begin{align}
  \min \quad & \sum_{v\in V} x_v \label{eq:ilp-obj}\\
  \text{s.t.} \quad & \sum_{u\in N_k(v)} x_u \;\ge\; 1, \quad v\in V
    \label{eq:ilp-cover}\\
  & x_v \in \{0,1\}, \quad v\in V. \label{eq:ilp-bin}
\end{align}
Both the number of binary variables and the number of constraints in
(\ref{eq:ilp-obj})--(\ref{eq:ilp-bin}) scale linearly with $n$, so
direct ILP solvers become impractical well before
$n=10^4$.

For $k=1$, sharper exponential-time analyses are available.
Grandoni~\cite{grandoni2006} reaches $\mathcal{O}(1.9053^n)$, while
van Rooij and Bodlaender~\cite{rooij2011} bring this down to
$\mathcal{O}(1.4969^n)$ in polynomial space.
Xiong and Xiao~\cite{xiong2024} couple LP-relaxation lower bounds with
a purely combinatorial branching rule and report roughly three orders
of magnitude of speed-up over the previous exact champion.
A complementary approach is the ILP-plus-enumeration solver of
Parra Inza, Vakhania, Sigarreta and Hern{\'a}ndez
Mira~\cite{parraInza2024EJOR}, which closes the optimality gap on a
substantial slice of medium-density inputs.

None of those exact lines of work survives the jump to social-network
scale. The decomposition strategy of Nguyen et~al.~\cite{nguyen2020} ---
partition $V$ via $k$-hop neighbourhoods, run a separate MILP on each
cluster through CPLEX, and stitch the partial solutions ---
remains the best published reference for million-vertex instances, at
the cost of thousands of seconds per instance and an external
commercial solver. A radius-bounded covering formulation underlies the
classical partition argument of Hassin and Tamir~\cite{hassin1991}.

\subsection{Greedy Heuristics}

Two greedy schedulers introduced by~\cite{nguyen2020} anchor the
construction stage we reuse. The first, denoted $\mathit{HEU_1}$ in
their paper, walks the vertex list in non-increasing degree order and
accepts each new vertex into $D$ only when it is still uncovered; we
call it \textbf{GD} (\emph{GreedyDegree}). The second, $\mathit{HEU_2}$,
is parameterised by an acceptance threshold $\theta\in\mathbb{Z}_{\ge 0}$
and weakens the admission test: a vertex is added either when it
itself remains uncovered or when its still-uncovered $k$-hop ball
contains at least $\theta$ vertices. Running the second scheduler
across a small grid $\Theta=\{0,1,\ldots,\theta_{\max}\}$ and keeping
the best returned $D$ gives the construction we call \textbf{GT}
(\emph{GreedyTheta}).

The conceptual difference matters. With $\theta\ge 1$, GT willingly
admits vertices whose own coverage requirement has already been met,
provided they contribute substantial extra coverage. The result is a
larger initial set in which many vertices are dominated by several
members of $D$ at once. This \emph{redundancy by design} is what the
post-optimisation phases of Section~\ref{sec:algo} subsequently
exploit: redundant or near-redundant dominators are exactly those that
RR removes and that LS replaces.

Two other baselines round out our construction set.
$\mathit{HEU_3}$ in~\cite{nguyen2020}, attributable to
Campan et~al.~\cite{campan2016}, is a partial-greedy variant tuned to
social graphs. Li, Potru and Shahrokhi~\cite{li2020} go further:
their LP-rounding heuristic reaches an integrality gap of $1.011$ on a
400K-vertex instance in $15$~seconds, and on graphs up to
7.7~million vertices, where the LP solver itself runs out of memory,
a hybrid construction is needed.

In our own pipeline we additionally report a $k$-core-aware ordering
(\textbf{KC}), an ordering by direct $|N_k(v)|$ size (\textbf{TH}), and
a HyperLogLog-based estimator of $|N_k(v)|$ (\textbf{SG}). The last is
a first attempt at the scalable techniques surveyed in
Section~\ref{sec:related:sketch}.

\subsection{Post-Optimization}

A redundant member of $D$ is one whose removal still leaves a feasible
solution. The straightforward way to find one is to delete each
candidate $v\in D$ tentatively and rebuild coverage on $N_k(v)$ from
scratch; the rebuild costs $O(|D|\,\Nk)$ for each $v$, and a single
sweep therefore reaches $O(|D|^2\,\Nk^2)$ in the worst case. The
post-optimisation block in $\mathit{HEU_4}$~\cite{nguyen2020}
operates in this spirit, with a coverage rebuild on every candidate
removal. On a graph with $|D|\!\approx\!3{,}300$ and $|N_3|\!\approx\!
3\!\times\!10^4$ such sweeps are observed to run for several minutes,
and frequently time out at higher $k$.

Recent work on plain MDS by Parra Inza, Sigarreta and
Vakhania~\cite{parraInza2024MDPI} reframes redundancy removal as a
\emph{purification} procedure driven by spanning-forest structure of
the greedy execution, and reports upper bounds up to seven times
sharper than earlier heuristics on more than $1{,}300$ benchmark
instances. Their analytical lens (spanning forests of the greedy
flowchart) and ours (per-vertex coverage counters) differ; we view our
linear-time procedure as the natural $k$-hop counterpart of their
$k=1$ result.

We are not aware of any prior work that achieves an $O(|D|\,\Nk)$
removal sweep for M$k$DSP through coverage-count bookkeeping, or that
extends the same primitive to a $(1,1)$-swap.

\subsection{Local Search and Metaheuristics}
\label{sec:related:ls}

For MDS, local search has been studied extensively. The dual-mode
local search DmDS of Zhu et~al.~\cite{zhu2024} interleaves $(2,1)$-
and $(3,2)$-swaps with frequency-based vertex selection and reports
state-of-the-art performance on common benchmarks. For the
\emph{weighted} variant, Chen et~al.'s
DeepOpt-MWDS~\cite{chen2023deepopt} introduces a multi-level
perturbation mechanism around a local-search core and outperforms
seven prior heuristics on massive-graph benchmarks. Abed
et~al.~\cite{abed2022hybrid} present a hybrid construction-plus-perturbation
local search.

For $k$-DS specifically, the only prior local-search work we are aware
of is RLS\_RCC by Li et~al.~\cite{li2022restart}: an Add/Swap/Drop
scheme combined with adaptive configuration checking and a restart
mechanism. RLS\_RCC establishes 40 new upper bounds across 54
general-graph instances and 32 new bounds on DIMACS instances, but
does not scale beyond medium-size graphs (no instance reported with
$n>10{,}000$). Our LS extends the local-search literature to $n=10^5$
at $k=3$ via a bounded $(1,1)$-swap that reuses the coverage-count
primitive of RR and is therefore amortized $O(|D|\cdot\Nk)$ per pass
with constant $L_{\max}$ and $C_{\max}$.

Population-based metaheuristics have also been applied to M$k$DSP:
Abed et~al.~\cite{abed2025} report a binary Equilibrium Optimizer on
small benchmarks, but their approach inherits the per-iteration cost
of population-based methods and does not scale to $10^5$ vertices.
General reviews of memetic, tabu, and ant-colony approaches to MDS
are surveyed in~\cite{haynes1998}.

\subsection{Sketch-Based and Streaming Methods}
\label{sec:related:sketch}

Estimating $|N_k(v)|$ at scale is a well-developed line of work in the
massive-graph community. Boldi, Rosa, and Vigna's
HyperANF~\cite{boldi2011} computes the neighborhood function for
graphs of billions of edges in hours via HyperLogLog counters and
broadword programming. Cohen's All-Distances Sketches with HIP
estimators~\cite{cohen2015} reduce variance by 50\% relative to prior
estimators and serve cardinality, distance-distribution, and
closeness queries. Chen et~al.~\cite{chen2023} give sublinear-space
streaming algorithms for estimating MDS \emph{size} (without producing
a dominating set) on sparse graphs.

None of these methods solves M$k$DSP directly, but they suggest a
clear path to scaling our pipeline to $10^9$ vertices: a HyperANF-style
preprocess can replace the $O(n\cdot\Nk)$ BFS pass at the heart of GT,
leaving the much smaller $O(|D|\cdot\Nk)$ work of RR and LS unchanged.
Our SG construction baseline is a first such hybrid; a fuller
exploration is left to future work.

\subsection{Learning-Based Approaches}

Graph neural networks have been applied to a range of combinatorial
problems on graphs, surveyed by Cappart
et~al.~\cite{cappart2023}. For an \emph{influence-maximization} variant
of $k$-domination, FastCover~\cite{ni2021} trains a graph
reversed-attention network in an unsupervised regime and reports a
$1000\times$ speedup over concurrent methods on graphs of up to
$400{,}000$ vertices. We emphasize, however, that FastCover solves the
\emph{budget-constrained} problem (maximize coverage given $|D|\le B$),
\emph{not} the minimum-cardinality M$k$DSP studied here. Chac{\'o}n
Sartori and Blum~\cite{chacon2022} similarly hybridize a biased
random-key genetic algorithm with a GNN scoring head for the same
budget-constrained variant.

Distributed approximation algorithms for MDS are studied by Alipour
and Salari~\cite{alipour2021}. None of these learning- or
distributed-system results currently provides a directly comparable
baseline for our pipeline.

\subsection{Preprocessing, Variants, and Applications}

Nguyen et~al.~\cite{nguyen2020} introduce a $k$-isolated-cluster (KIC)
preprocessing step that contracts dense neighborhoods. While effective
for $k=1,2$, our experiments show that KIC provides no benefit for
$k\ge 3$ on the benchmark studied here
(Section~\ref{sec:exp:heuristics}).

M$k$DSP extends naturally to weighted
variants~\cite{barrena2025} and to digraphs and road
networks~\cite{dijkstra2024}. Foundational applications of
$k$-domination include max-min cluster formation in wireless ad-hoc
networks~\cite{amis2000} and influence maximization in
social networks~\cite{kempe2003}. Hardness and inapproximability of
the underlying covering problem are due to Garey and
Johnson~\cite{garey1979} and Raz and Safra~\cite{raz1997}.


% =============================================================
\section{Proposed Algorithms}
\label{sec:algo}
% =============================================================

\subsection{GreedyTheta (GT) — Construction Phase}

GreedyTheta sweeps the vertex list in descending degree order $t$
times, once per threshold in a small finite set
$\Theta=\{\theta_1,\ldots,\theta_t\}\subset\mathbb{Z}_{\ge 0}$, and
keeps the smallest dominating set encountered. The sweep for a
particular $\theta$ admits $v$ into $D$ whenever $v$ is not already
covered or whenever the still-uncovered slice of $N_k(v)$ has
cardinality at least $\theta$. Pseudocode is given as
Algorithm~\ref{alg:gt}.

\begin{algorithm}[ht]
\caption{GreedyTheta$(G, k, \Theta)$}
\label{alg:gt}
\begin{algorithmic}[1]
\Require Graph $G=(V,E)$, integer $k$, threshold set $\Theta$
\Ensure A valid $k$-dominating set $D\subseteq V$
\State $D^*\gets\emptyset$;\; $\text{best}\gets\infty$
\State $\text{order}\gets\text{sort } V \text{ by } \deg(v) \text{ descending}$
\For{each $\theta\in\Theta$}
  \State $D\gets\emptyset$;\; $\text{covered}\gets\emptyset$
  \For{each $v\in\text{order}$}
    \State $N_k\gets N_k(v)$ \Comment{BFS $k$ hops from $v$}
    \If{$v\notin\text{covered}$ \textbf{or} $|N_k\setminus\text{covered}|\ge\theta$}
      \State $D\gets D\cup\{v\}$
      \State $\text{covered}\gets\text{covered}\cup N_k$
    \EndIf
  \EndFor
  \If{$|D|<\text{best}$}
    \State $D^*\gets D$;\; $\text{best}\gets|D|$
  \EndIf
\EndFor
\State \Return $D^*$
\end{algorithmic}
\end{algorithm}

\textbf{Role of $\theta$.}
At $\theta=0$ the admission rule reduces to the conservative one of
GD: only uncovered candidates ever enter $D$. Each additional unit of
$\theta$ relaxes the rule a step further, letting GT pull in already
covered vertices when they bring an extra $\theta$ vertices into the
coverage zone. The set returned at $\theta\ge 1$ is therefore
deliberately larger than the GD output, and—as we exploit
immediately below—its redundant fraction is also larger.

\textbf{Complexity.} A single BFS to depth $k$ touches $\Nk$
vertices on average, line~6 runs $n$ times per $\theta$, and there
are $|\Theta|$ outer iterations; the total cost is
$O(|\Theta|\cdot n\cdot\Nk)$.

\subsection{Fast Redundant Removal (RR)}
\label{sec:rr}

Starting from any feasible $D$, RR sweeps $D$ once and discards every
member whose distance-$k$ ball is already doubly covered. The
expensive part of the classical procedure—rebuilding the coverage
status from scratch after each tentative removal—is avoided by
maintaining the per-vertex counter $\cov(\cdot)$ explicitly and
updating it in place. Pseudocode follows.

\begin{algorithm}[ht]
\caption{FastRedundantRemoval$(G, k, D)$}
\label{alg:rr}
\begin{algorithmic}[1]
\Require Graph $G=(V,E)$, integer $k$, valid $k$-dominating set $D$
\Ensure $D'\subseteq D$, still a valid $k$-dominating set, $|D'|\le|D|$
\Statex \textit{// Phase 1: precompute $k$-hop neighborhoods for $D$}
\For{each $v\in D$}
  \State $\text{nk}[v]\gets N_k(v)$ \Comment{BFS $k$ hops}
\EndFor
\Statex \textit{// Phase 2: compute coverage counts}
\For{each $v\in D$}
  \For{each $u\in\text{nk}[v]$}
    \State $\cov[u]\gets\cov[u]+1$
  \EndFor
\EndFor
\Statex \textit{// Phase 3: sort by neighborhood size ascending}
\State $\text{sorted\_D}\gets\text{sort } D \text{ by } |\text{nk}[v]| \text{ ascending}$
\Statex \textit{// Phase 4: greedy removal}
\State $\text{removed}\gets\emptyset$
\For{each $v\in\text{sorted\_D}$}
  \If{$\min\{\cov[u]:u\in\text{nk}[v]\}\ge 2$}
    \State $\text{removed}\gets\text{removed}\cup\{v\}$
    \For{each $u\in\text{nk}[v]$}
      \State $\cov[u]\gets\cov[u]-1$
    \EndFor
  \EndIf
\EndFor
\State \Return $D\setminus\text{removed}$
\end{algorithmic}
\end{algorithm}

\textbf{Key design decisions.}
\begin{enumerate}[nosep,label=(\alph*)]
\item \textbf{Ascending-size ordering (line~10).}
  Vertices with smaller $k$-hop neighborhoods are considered first.
  Removing such a vertex decreases coverage counts for fewer nodes,
  preserving more removal opportunities for vertices considered later.
\item \textbf{Coverage-count check (line~12).}
  Instead of rebuilding covered sets, we only check that every node in
  $N_k(v)$ has $\cov[u]\ge 2$. If so, removing $v$ still leaves
  $\cov[u]\ge 1$ everywhere.
\item \textbf{Decremental update (lines~14--15).}
  After removing $v$, we decrement $\cov[u]$ for all $u\in N_k(v)$ in
  $O(|N_k(v)|)$ time.
\end{enumerate}

\begin{theorem}[Correctness of FastRedundantRemoval]
\label{thm:rr-correct}
Algorithm~\ref{alg:rr} returns a valid $k$-dominating set.
\end{theorem}

\begin{proof}
We maintain the following invariant.

\textbf{Invariant~(I).} For all $u\in V$, $\cov[u]\ge 1$ after each
iteration of the loop on line~11.

\textbf{Base case.} Before any removal, $D$ is a valid $k$-dominating
set, so $\cov[u]\ge 1$ for every $u\in V$.

\textbf{Inductive step.} Suppose~(I) holds before considering vertex
$v\in D$. If $v$ is selected for removal (line~13), then
$\cov[u]\ge 2$ for every $u\in\text{nk}[v]$ by the test on line~12.
After decrementing (line~15), $\cov[u]\ge 1$ for every $u\in\text{nk}[v]$.
For all $u\notin\text{nk}[v]$, $\cov[u]$ is unchanged.
Hence~(I) is preserved. If $v$ is not removed, no coverage count
changes and~(I) holds trivially.

After processing every vertex,~(I) ensures $\cov[u]\ge 1$ for every
$u\in V$, so $D'=D\setminus\text{removed}$ is a valid $k$-dominating set.
\end{proof}

\begin{theorem}[Complexity of FastRedundantRemoval]
\label{thm:rr-complexity}
The time complexity of Algorithm~\ref{alg:rr} is $O(|D|\cdot\Nk)$
where $\Nk = \frac{1}{|D|}\sum_{v\in D}|N_k(v)|$.
\end{theorem}

\begin{proof}
\begin{itemize}[nosep]
  \item \textbf{Phase~1} (BFS): $\sum_{v\in D}O(|N_k(v)|)=O(|D|\cdot\Nk)$.
  \item \textbf{Phase~2} (coverage counts): same sum, $O(|D|\cdot\Nk)$.
  \item \textbf{Phase~3} (sorting):
    $O(|D|\log|D|)\le O(|D|\cdot\Nk)$ when $\Nk\ge\log|D|$, which holds
    on graphs of practical interest.
  \item \textbf{Phase~4} (removal): for each $v\in D$, the minimum check
    and possible decrement cost $O(|N_k(v)|)$, totaling $O(|D|\cdot\Nk)$.
\end{itemize}
Overall: $O(|D|\cdot\Nk)$.
\end{proof}

Table~\ref{tab:complexity} summarizes the comparison with classical
removal.

\begin{table}[ht]
\centering
\caption{Complexity: classical vs.\ fast redundant removal.}
\label{tab:complexity}
\begin{tabular}{lcc}
\toprule
& \textbf{Classical RR} & \textbf{Fast RR (ours)}\\
\midrule
Per vertex & $O(|D|\cdot\Nk)$ & $O(\Nk)$\\
Total      & $O(|D|^2\cdot\Nk^2)$ & $O(|D|\cdot\Nk)$\\
\bottomrule
\end{tabular}
\end{table}


\subsection{Local Swap (LS)}
\label{sec:ls}

After RR has converged, the dominating set may still hold vertices
whose loss set is small but not empty: they are not strictly
redundant, yet a single outside vertex could absorb their unique
contribution. LS scans $D$ in passes, removing the strictly redundant
members (Case~A below, identical to one step of RR) and otherwise
attempting a $(1,1)$-exchange with such an outside substitute
(Case~B).

\begin{algorithm}[ht]
\caption{LocalSwap$(G,k,D,P,L_{\max},C_{\max})$}
\label{alg:ls}
\begin{algorithmic}[1]
\Require Graph $G$, integer $k$, valid set $D$, max passes $P$, loss
  bound $L_{\max}$, candidate bound $C_{\max}$
\Ensure $D'$, valid $k$-dominating set with $|D'|\le|D|$
\For{$\text{pass}=1,\ldots,P$}
  \State Recompute $\text{nk}[v]\gets N_k(v)$ for $v\in D$
  \State Recompute $\cov[u]$ from scratch
  \State $\text{sorted\_D}\gets\text{sort } D$ by $|\text{nk}[v]|$ ascending
  \State $\text{improved}\gets\textsc{false}$
  \For{each $v\in\text{sorted\_D}$ with $v\in D$}
    \State $\loss\gets\{u\in\text{nk}[v]:\cov[u]\le 1\}$
    \If{$\loss=\emptyset$} \Comment{Case~A: $v$ redundant}
      \State $D\gets D\setminus\{v\}$;\; decrement $\cov$ on $\text{nk}[v]$
      \State $\text{improved}\gets\textsc{true}$;\; \textbf{continue}
    \EndIf
    \If{$|\loss|>L_{\max}$} \textbf{continue} \EndIf
    \Statex \textit{// Case~B: collect candidates near $\loss$}
    \State $\text{cand}\gets\bigcup_{u\in\loss}\bigl(N_k(u)\setminus D\bigr)$,
       capped at $2C_{\max}$
    \State $\text{cand\_list}\gets$ first $C_{\max}$ of $\text{cand}$ sorted by
       $\deg(\cdot)$ descending
    \State $w^*\gets\textsc{none}$;\; $\text{best}\gets-1$
    \For{$w\in\text{cand\_list}$}
      \State $\text{nk}_w\gets N_k(w)$
      \If{$\loss\subseteq\text{nk}_w$}
        \State $\text{extra}\gets|\{u\in\text{nk}_w\setminus\loss:\cov[u]\le 1\}|$
        \If{$\text{extra}>\text{best}$}
          \State $w^*\gets w$;\; $\text{nk}_{w^*}\gets\text{nk}_w$;\; $\text{best}\gets\text{extra}$
        \EndIf
      \EndIf
    \EndFor
    \If{$w^*\ne\textsc{none}$}
      \State $D\gets(D\setminus\{v\})\cup\{w^*\}$
      \State decrement $\cov$ on $\text{nk}[v]$;\; increment $\cov$ on $\text{nk}_{w^*}$
      \State $\text{nk}[w^*]\gets\text{nk}_{w^*}$;\; $\text{improved}\gets\textsc{true}$
    \EndIf
  \EndFor
  \If{$\neg\text{improved}$} \textbf{break} \EndIf
\EndFor
\State \Return $D$
\end{algorithmic}
\end{algorithm}

\textbf{Design notes.}
The algorithm interleaves the redundant-removal logic of
Algorithm~\ref{alg:rr} (Case~A) with a $(1,1)$-swap (Case~B). The
candidate set in line~12 is the union of $k$-hop neighborhoods of
$\loss$, capped to $2C_{\max}$ vertices to bound the candidate-collection
work. The tie-breaking rule in lines~17--19 prefers a substitute $w$
that brings the most additional uniquely covered vertices into a
``shielded'' state with $\cov\ge 2$, opening further removal opportunities
in subsequent passes.

\begin{lemma}[Feasibility under LocalSwap]
\label{lem:ls-correct}
Algorithm~\ref{alg:ls} returns a valid $k$-dominating set after every
pass.
\end{lemma}

\begin{proof}
We show the invariant $\cov[u]\ge 1$ for all $u\in V$ is preserved by
each modification of $D$ inside the inner loop.

\textbf{Case~A (line~9 -- removal).} Identical to the inductive step in
Theorem~\ref{thm:rr-correct}: the test $\loss=\emptyset$ on line~8 means
$\cov[u]\ge 2$ for every $u\in\text{nk}[v]$, so decrementing leaves
$\cov[u]\ge 1$ on $\text{nk}[v]$ and unchanged elsewhere.

\textbf{Case~B (lines~21--23 -- swap).} Before the swap, the test on
line~16 guarantees $\loss\subseteq\text{nk}_{w^*}$. After decrementing
$\cov$ on $\text{nk}[v]$, every $u\in\text{nk}[v]\setminus\loss$ had
$\cov[u]\ge 2$ and now has $\cov[u]\ge 1$, while every $u\in\loss$ now
has $\cov[u]=0$. The subsequent increment on $\text{nk}_{w^*}$ raises
$\cov[u]$ by one for every $u\in\loss$ (since $\loss\subseteq\text{nk}_{w^*}$),
restoring $\cov[u]\ge 1$. Vertices outside $\text{nk}[v]\cup\text{nk}_{w^*}$
are unchanged.

The invariant is therefore preserved by every modification, and at the
end of each pass $D$ is a valid $k$-dominating set.
\end{proof}

\begin{theorem}[Complexity of LocalSwap]
\label{thm:ls-complexity}
With $L_{\max}$ and $C_{\max}$ treated as constants, Algorithm~\ref{alg:ls}
runs in $O(P\cdot|D|\cdot\Nk)$ time.
\end{theorem}

\begin{proof}
A single pass costs (i) $O(|D|\cdot\Nk)$ for line~2 (BFS),
(ii) $O(|D|\cdot\Nk)$ for line~3 (coverage), (iii) $O(|D|\log|D|)$ for
line~4 (sort), and (iv) for each vertex $v$, an $O(|N_k(v)|)$
loss-set computation plus, if Case~B fires, an
$O(L_{\max}\cdot\Nk)$ candidate-collection (line~12) and an
$O(C_{\max}\cdot\Nk)$ candidate-evaluation (lines~14--19). With
$L_{\max},C_{\max}=O(1)$, the per-vertex cost is $O(\Nk)$ and the
inner loop totals $O(|D|\cdot\Nk)$. Summing across $P$ passes yields
$O(P\cdot|D|\cdot\Nk)$.
\end{proof}

\subsection{Pipelines}
\label{sec:pipelines}

We compose the building blocks into two pipelines.

\begin{description}[nosep,labelwidth=2.6cm,leftmargin=2.8cm]
  \item[GT$+$RR (2 phases):]
    Run GT with $\Theta=\{0,1,2,3,4\}$, then RR.
    Bound $O(|\Theta|\cdot n\cdot\Nk)+O(|D|\cdot\Nk)$.
  \item[GT$+$LS$+$RR (3 phases):]
    Run GT, then LS with $P=10,\,L_{\max}=5,\,C_{\max}=200$, then RR
    on the LS output.
    Bound $O(|\Theta|\cdot n\cdot\Nk)+O(P\cdot|D|\cdot\Nk)$.
\end{description}

\textbf{Why LS before RR?}
When LS performs a swap, the substitute $w^*$ is chosen to maximize
``extra'' uniquely covered vertices that become doubly covered
(line~17). Subsequent passes therefore find more redundant vertices.
Running RR \emph{after} LS lets RR exploit this extra redundancy. We
confirm experimentally in Section~\ref{sec:exp:main} that LS$+$RR
outperforms RR$+$LS by 3--10 vertices on average per graph.

No preprocessing is required for $k\ge 3$. KIC preprocessing of
\cite{nguyen2020} can reduce graph size for $k=1,2$ but provides no
benefit at $k=3$ on this benchmark
(Section~\ref{sec:exp:heuristics}).


% =============================================================
\section{Experimental Evaluation}
\label{sec:experiments}
% =============================================================

\subsection{Setup}
\label{sec:exp:setup}

\textbf{Implementation.} Python~3.11 with NetworkX~3.4.2, single thread,
no GPU. Hardware: Intel Core i7-12700H CPU, 32~GB RAM, Ubuntu 22.04
LTS. All wall-clock times are reported as measured.

\textbf{Benchmark.} We use 9~Erd\H{o}s--R\'enyi-like random graphs
generated by the \texttt{igraph} package, each with $n=100{,}000$
vertices.
Their edge counts and average degrees are summarized in
Table~\ref{tab:graph-stats} (extracted directly from the GraphML files).
We refer to them as \texttt{rand\_100k\_1}--\texttt{rand\_100k\_9}.
We additionally use one synthetic graph
\texttt{graphbe} ($n=614{,}548$, $|E|=630{,}469$,
average degree~$2.05$ — sparse but much larger) as a stress-test instance
(Section~\ref{sec:exp:graphbe}).

\begin{table}[ht]
\centering
\caption{Benchmark graph statistics ($n=100{,}000$, undirected).
  Edge counts read from the GraphML headers.}
\label{tab:graph-stats}
\begin{tabular}{lrrr}
\toprule
\textbf{Graph} & $|V|$ & $|E|$ & \textbf{avg.\ deg.} \\
\midrule
rand\_100k\_1 & 100{,}000 & 249{,}725 & 4.99 \\
rand\_100k\_2 & 100{,}000 & 249{,}893 & 5.00 \\
rand\_100k\_3 & 100{,}000 & 249{,}991 & 5.00 \\
rand\_100k\_4 & 100{,}000 & 249{,}681 & 4.99 \\
rand\_100k\_5 & 100{,}000 & 249{,}823 & 5.00 \\
rand\_100k\_6 & 100{,}000 & 249{,}528 & 4.99 \\
rand\_100k\_7 & 100{,}000 & 249{,}867 & 5.00 \\
rand\_100k\_8 & 100{,}000 & 249{,}640 & 4.99 \\
rand\_100k\_9 & 100{,}000 & 249{,}577 & 4.99 \\
\bottomrule
\end{tabular}
\end{table}

\textbf{Parameter settings.}
\begin{itemize}[nosep]
  \item GT: $\Theta=\{0,1,2,3,4\}$, time limit 60\,s per $\theta$ (rarely
    triggered at $n=10^5$).
  \item RR: no parameters (fully automatic).
  \item LS: $P=10$ passes, $L_{\max}=5$, $C_{\max}=200$. These are the
    defaults of our open-source implementation; we did not tune them.
  \item $k\in\{1,2,3\}$, with $k=3$ as the primary setting.
\end{itemize}

\textbf{Construction baselines for comparison.}
GD (GreedyDegree), KC (KCoreGreedy), TH (TwoHopGreedy), and SG (sketch
greedy with HyperLogLog) accompany GT in
Section~\ref{sec:exp:heuristics}.

\textbf{Validity check.} Every reported $|D|$ corresponds to a verified
valid $k$-dominating set: the implementation rebuilds coverage from
scratch on the final solution and confirms $\cov(u)\ge 1$ for every
$u\in V$.


\subsection{Main Results: Pipelines on the 100K Benchmark}
\label{sec:exp:main}

Table~\ref{tab:main} reports $|D|$ for the five GT-based configurations
across all 9~graphs at $k=3$. Every cell is a single measured run on
the corresponding instance; averages on the last row are unweighted.

\begin{table}[ht]
\centering
\caption{$|D|$ on rand\_100k\_$\{1,\ldots,9\}$, $k=3$.
  All configurations validated.}
\label{tab:main}
\small
\begin{tabular}{lrrrrr}
\toprule
\textbf{Graph}
  & \textbf{GT}
  & \textbf{GT$+$RR}
  & \textbf{GT$+$RR$+$LS}
  & \textbf{GT$+$LS}
  & \textbf{GT$+$LS$+$RR}\\
\midrule
rand\_100k\_1 & 3{,}649 & 3{,}257 & 2{,}548 & 2{,}551 & 2{,}549 \\
rand\_100k\_2 & 3{,}700 & 3{,}303 & 2{,}580 & 2{,}575 & 2{,}570 \\
rand\_100k\_3 & 3{,}661 & 3{,}249 & 2{,}552 & 2{,}558 & 2{,}554 \\
rand\_100k\_4 & 3{,}710 & 3{,}319 & 2{,}590 & 2{,}590 & 2{,}587 \\
rand\_100k\_5 & 3{,}613 & 3{,}228 & 2{,}529 & 2{,}530 & 2{,}529 \\
rand\_100k\_6 & 3{,}752 & 3{,}364 & 2{,}592 & 2{,}596 & 2{,}593 \\
rand\_100k\_7 & 3{,}685 & 3{,}314 & 2{,}626 & 2{,}617 & 2{,}613 \\
rand\_100k\_8 & 3{,}648 & 3{,}277 & 2{,}566 & 2{,}593 & 2{,}585 \\
rand\_100k\_9 & 3{,}758 & 3{,}342 & 2{,}624 & 2{,}621 & 2{,}614 \\
\midrule
\textbf{Average}
  & \textbf{3{,}686} & \textbf{3{,}295} & \textbf{2{,}579}
  & \textbf{2{,}581} & \textbf{2{,}577}\\
\textbf{vs.\ GT}
  & --- & $-10.6\%$ & $-30.0\%$ & $-29.9\%$ & $-30.1\%$\\
\textbf{vs.\ GT$+$RR}
  & $+11.9\%$ & --- & $-21.7\%$ & $-21.7\%$ & $-21.8\%$\\
\bottomrule
\end{tabular}
\end{table}

\textbf{Findings.}
\begin{itemize}[nosep]
  \item \textbf{RR alone gives a consistent ${\sim}10.6\%$ reduction.}
    Across the 9~graphs, $\Delta_{\text{RR}}$ ranges from $-10.1\%$
    (rand\_100k\_7) to $-11.3\%$ (rand\_100k\_3); the average is
    $(3{,}686-3{,}295)/3{,}686=10.6\%$.
  \item \textbf{LS dominates as the strongest improver.}
    Adding LS atop GT$+$RR drops $|D|$ by a further $21.8\%$, from
    $3{,}295$ to $2{,}577$.
  \item \textbf{LS$+$RR ordering wins by a small but consistent margin.}
    GT$+$LS$+$RR achieves $2{,}577$ vs.\ $2{,}579$ for GT$+$RR$+$LS,
    confirming the heuristic argument of
    Section~\ref{sec:pipelines}: running LS first creates additional
    swap-induced redundancy that RR can clean up.
\end{itemize}


\subsection{Time Decomposition}
\label{sec:exp:time}

Table~\ref{tab:time} reports per-phase running time, averaged over the
9~graphs at $k=3$.

\begin{table}[ht]
\centering
\caption{Wall-clock time (seconds), average over 9 graphs, $k=3$.}
\label{tab:time}
\begin{tabular}{lrrr}
\toprule
\textbf{Pipeline} & \textbf{GT phase} & \textbf{Post-opt phase} & \textbf{Total}\\
\midrule
GT             & ${\sim}30$ & ---           & ${\sim}30$ \\
GT$+$RR        & ${\sim}30$ & ${\sim}1$     & ${\sim}31$ \\
GT$+$RR$+$LS   & ${\sim}30$ & ${\sim}154$   & ${\sim}184$ \\
GT$+$LS$+$RR   & ${\sim}30$ & ${\sim}156$   & ${\sim}186$ \\
\bottomrule
\end{tabular}
\end{table}

The RR phase by itself adds at most $0.2$--$2.6$\,s of overhead on the
9~graphs (verified per-graph in
Appendix~\ref{app:per-graph-time}), confirming the $O(|D|\cdot\Nk)$
bound of Theorem~\ref{thm:rr-complexity} is realized in practice. LS is
two orders of magnitude more expensive but contributes the bulk of the
quality improvement.


\subsection{Comparison of Construction Heuristics}
\label{sec:exp:heuristics}

Table~\ref{tab:heuristics} compares the four construction baselines and
their RR variants on \texttt{rand\_100k\_2} ($k=3$).

\begin{table}[ht]
\centering
\caption{Heuristic comparison on \texttt{rand\_100k\_2}, $k=3$.}
\label{tab:heuristics}
\begin{tabular}{lrrr}
\toprule
\textbf{Pipeline} & $|D|$ & \textbf{Valid} & \textbf{Time (s)}\\
\midrule
GD                 & 4{,}921          & \checkmark & 0.30 \\
GT                 & 3{,}700          & \checkmark & 40.25 \\
KC                 & 4{,}926          & \checkmark & 2.14 \\
TH                 & 4{,}932          & \checkmark & 7.98 \\
\midrule
GD$+$RR            & 4{,}921          & \checkmark & 1.14 \\
\textbf{GT$+$RR}   & \textbf{3{,}303} & \checkmark & 33.97 \\
KC$+$RR            & 4{,}926          & \checkmark & 1.75 \\
TH$+$RR            & 4{,}932          & \checkmark & 7.35 \\
\midrule
KIC$+$GD$+$RR      & 4{,}921          & \checkmark & 44.24 \\
KIC$+$GT$+$RR      & 3{,}303          & \checkmark & 64.72 \\
KIC$+$KC$+$RR      & 4{,}926          & \checkmark & 36.44 \\
KIC$+$TH$+$RR      & 4{,}932          & \checkmark & 41.32 \\
\bottomrule
\end{tabular}
\end{table}

\textbf{Findings.}
\begin{enumerate}[nosep]
  \item \textbf{GT dominates the other constructions} by 24--33\%
    ($3{,}700$ vs.\ $4{,}921$/$4{,}926$/$4{,}932$).
  \item \textbf{RR only reduces GT solutions; GD/KC/TH are
    near-minimal already.} For GD/KC/TH the test
    $\min\cov\ge 2$ is never satisfied, so no vertex is removed
    ($|D|$ unchanged after RR). GT's threshold mechanism is what
    creates the redundancy that RR exploits.
  \item \textbf{KIC adds ${\sim}30$\,s with no benefit at $k=3$.}
    KIC$+$GT$+$RR returns the same $|D|=3{,}303$ as GT$+$RR but takes
    $64.72$\,s versus $33.97$\,s. We therefore drop KIC for $k=3$ in
    the main pipelines of Section~\ref{sec:exp:main}.
\end{enumerate}

LS atop the inferior constructions further demonstrates the
``LS erases construction quality gaps'' observation
(Section~\ref{sec:discussion:why-ls}). Across the 9~graphs, the average
$|D|$ for the GD$+$LS$+$RR pipeline is $2{,}584$, only $0.27\%$ worse
than GT$+$LS$+$RR's $2{,}577$, despite GD's standalone average being
$33.6\%$ worse than GT's
($4{,}923$ vs.\ $3{,}686$).
SG$+$LS$+$RR achieves $2{,}582$ ($0.19\%$ worse than GT$+$LS$+$RR).


\subsection{Impact of $k$}
\label{sec:exp:k}

Table~\ref{tab:impact-k} reports $|D|$ as $k$ varies on
\texttt{rand\_100k\_2}.

\begin{table}[ht]
\centering
\caption{Impact of $k$ on \texttt{rand\_100k\_2}.}
\label{tab:impact-k}
\begin{tabular}{crrrrl}
\toprule
$k$ & \textbf{GD} & \textbf{GT} & \textbf{GT$+$RR}
  & $\Delta_{\text{RR}}$ & \textbf{RR effective?}\\
\midrule
1 & 30{,}741 & 28{,}441 & 28{,}441 & $0.0\%$ & No \\
2 & 10{,}634 &  9{,}853 &  9{,}853 & $0.0\%$ & No \\
3 &  4{,}921 &  3{,}700 &  3{,}303 & $-10.7\%$ & \textbf{Yes}\\
\bottomrule
\end{tabular}
\end{table}

For small $k$, the $k$-hop neighborhoods are small (avg.\ $|N_1|\approx 6$,
$|N_2|\approx 30$ on this graph), GT does not over-cover, and no
redundancy emerges. For $k=3$ (avg.\ $|N_3|\approx 150$), neighborhoods
overlap significantly, GT exploits the $\theta$-mechanism, and RR
removes ${\sim}10.7\%$ of the dominating set. Fast RR is therefore most
beneficial precisely when the graph's $k$-hop structure produces a
high-redundancy initial solution.


\subsection{Case Study: \texttt{graphbe}}
\label{sec:exp:graphbe}

A complementary stress test on a sparse but much larger synthetic graph
(\texttt{graphbe}, $n=614{,}548$, $|E|=630{,}469$,
average degree $2.05$) yields the results in
Table~\ref{tab:graphbe}.

\begin{table}[ht]
\centering
\caption{\texttt{graphbe}, $k=3$. LS attempts no removals and
  performs only $40$ swaps over $10$ passes, yielding no further size
  reduction.}
\label{tab:graphbe}
\begin{tabular}{lrr}
\toprule
\textbf{Pipeline} & $|D|$ & \textbf{Time (s)}\\
\midrule
GT             & 497 & 300.77 \\
GT$+$RR        & 468 & 307.68 \\
GT$+$RR$+$LS   & 468 & 341.43 \\
\bottomrule
\end{tabular}
\end{table}

On this instance RR removes $29$~vertices (a 5.8\% reduction) but LS
finds no useful $(1,1)$-swap: every candidate $v\in D$ either has
$|\loss(v)|>L_{\max}$ or no candidate $w$ within the bounded search
covers its loss. The instance illustrates an LS \emph{floor} regime
where RR has already reached a local optimum that LS cannot escape with
$P=10,\,L_{\max}=5,\,C_{\max}=200$.


\subsection{Comparison with Prior Work}
\label{sec:exp:vs-prior}

A direct head-to-head comparison with the $\mathit{HEU_4}$ pipeline of
Nguyen et~al.~\cite{nguyen2020} on the rand\_100k benchmark would
require re-running their CPLEX-based partition MILP on the same
machine; we do not have such matched-instance data and refrain from a
numeric speedup claim. For context, their published Table~4 reports
$\mathit{HEU_4}$ running times of $1{,}695$--$42{,}600$~seconds
($k=3$) on real-world social-network instances of comparable scale
(\texttt{soc-delicious}, $|V|=536$K to \texttt{soc-livejournal},
$|V|=4{,}033$K), all relying on a commercial MILP solver. Our pipelines
finish in $30$\,s (GT$+$RR) to $186$\,s (GT$+$LS$+$RR) on a single
thread without any MILP solver. A controlled comparison on identical
instances is left to future work.


% =============================================================
\section{Discussion}
\label{sec:discussion}
% =============================================================

\subsection{Why GT Creates Beneficial Redundancy}

The role of $\theta$ in GT can be read off
Table~\ref{tab:heuristics} directly. The three baselines whose
acceptance rule never admits an already-covered vertex (GD, KC, TH)
yield dominating sets that RR cannot shrink at all: the rule by
construction produces no doubly-covered vertices, so the minimum on
line~\ref{alg:rr}, line~12 is always one. GT with $\theta\ge 1$, on
the contrary, deliberately accepts vertices whose contribution is
\emph{additional} coverage rather than first-time coverage, and the
result is a set in which $\cov(u)\ge 2$ for many $u$. RR is then no
longer ineffectual: the same line~12 fires often enough to remove
about a tenth of $D$ at $k=3$.

Phrased differently, GT trades an immediate increase in $|D|$ for a
gain in \emph{redundancy mass} downstream, which RR turns back into a
net decrease in $|D|$. Without $\theta$, the slack RR feeds on does
not exist.


\subsection{Why LocalSwap Dominates RR Alone}
\label{sec:discussion:why-ls}

LS does not merely repeat the work of RR. Whenever
$\loss(v)$ has size between one and $L_{\max}$, RR is structurally
unable to remove $v$, while LS may still replace $v$ with a single
substitute. The difference is structural rather than tuning-related:
even at $L_{\max}=5,\,C_{\max}=200,\,P=10$, all comfortable defaults,
LS already accounts for the additional $20.5$ percentage points of
reduction reported in Table~\ref{tab:main}.

A consequence visible in our experiments is that the choice of
construction loses importance once LS is in the pipeline. GD, SG and
GT differ by up to a third on the standalone column of
Table~\ref{tab:main}, but their three corresponding LS$+$RR runs end
within roughly $0.3\%$ of one another. The construction phase
controls the runtime profile (a fast SG is still useful), but no
longer the asymptotic solution quality.


\subsection{Coverage Count as a Cross-Procedure Primitive}

A single per-vertex counter, decremental updates of cost
$O(|N_k(v)|)$, and three predicates suffice to express everything in
RR and LS:
\begin{itemize}[nosep]
  \item the redundancy test $\min_{u\in N_k(v)}\cov(u)\ge 2$;
  \item the loss set $\loss(v)=\{u\in N_k(v):\cov(u)\le 1\}$;
  \item the swap admissibility check $\loss(v)\subseteq N_k(w)$.
\end{itemize}
The same recipe transfers without modification to other covering
problems with a radius parameter—plain set cover, distance-bounded
facility location, $r$-dominating set—where redundancy removal and
small-radius swaps are part of the standard tool-kit.


\subsection{Limitations}

\begin{enumerate}[nosep]
  \item \textbf{Graph family.} The 100K benchmark consists exclusively
    of Erd\H{o}s--R\'enyi-like random graphs with average degree~$5$.
    Power-law and small-world networks of comparable scale may behave
    differently, both for RR (which depends on $\theta$-induced
    redundancy) and for LS (which depends on the existence of
    high-degree substitutes near every $\loss(v)$). We did not test on
    real-world social networks.
  \item \textbf{$k$ regime.} RR is empirically ineffective for
    $k=1,2$ on this benchmark
    (Table~\ref{tab:impact-k}); the gain reported in
    Section~\ref{sec:exp:main} is specific to $k=3$.
  \item \textbf{Optimality gap.} Without a lower bound, we cannot bound
    the approximation ratio. LP relaxation
    (\ref{eq:ilp-cover})--(\ref{eq:ilp-bin}) provides such a bound but
    its computation at $n=10^5$ is itself nontrivial and is left to
    future work.
  \item \textbf{Bounded LS search.} With $L_{\max}=5$ and
    $C_{\max}=200$, LS may overlook valid swaps when the loss set is
    larger or when no high-degree substitute is among the first
    $C_{\max}$ candidates. We did not tune $L_{\max}$ or $C_{\max}$.
  \item \textbf{No cross-machine timing comparison.} The MILP-based
    timings of~\cite{nguyen2020} are not directly comparable to ours
    (different machines, different instances). We therefore make no
    speedup claim.
\end{enumerate}


% =============================================================
\section{Conclusion}
\label{sec:conclusion}
% =============================================================

Two short procedures, both organised around a single coverage
counter, were studied here. Fast Redundant Removal sweeps $D$ once
and erases every member whose distance-$k$ ball is already
double-dominated; the sweep runs in $O(|D|\cdot\Nk)$ time, an
improvement from the textbook $O(|D|^2\cdot\Nk^2)$. Local Swap pushes
the same counter past the redundancy threshold, tolerating a bounded
amount of uniquely-covered ``loss'' in exchange for permission to
swap the offending dominator with a single outside vertex; with the
bounds $L_{\max},C_{\max}$ treated as constants its cost is
$O(P\cdot|D|\cdot\Nk)$ over $P$ passes. Both procedures keep the
coverage counter at least one everywhere, and so never produce an
infeasible intermediate set.

On the nine fresh random graphs with $n=10^5$ used here, GT$+$RR
shrinks the GreedyTheta construction by $10.6\%$ in roughly one
extra second per instance (worst case four seconds), and the full
three-phase pipeline GT$+$LS$+$RR shrinks it by $30.1\%$ in roughly
$185$~seconds. Effectively all of the improvement beyond RR comes
from LS; the differences between greedy-degree, sketch-based and
$\theta$-thresholded constructions are absorbed by the local-search
stage to within a fraction of a percent.

\textbf{Future directions.}
Five threads suggest themselves naturally.
(i)~A real-world counterpart of our benchmark — power-law and
small-world graphs, ideally the same instances on which Nguyen
et~al.~\cite{nguyen2020} reported their MILP-based results — would
allow a matched head-to-head comparison with $\mathit{HEU_4}$ that we
deliberately avoided here for lack of data.
(ii)~Replacing the $O(n\cdot\Nk)$ BFS of GT with a HyperLogLog-style
estimator of $|N_k(v)|$ as in HyperANF~\cite{boldi2011} would let
the construction stage scale to $10^6$--$10^9$ vertices while leaving
the $O(|D|\cdot\Nk)$ RR/LS stages untouched.
(iii)~An LP lower bound on (\ref{eq:ilp-obj})--(\ref{eq:ilp-bin})
would let us bound the optimality gap, where right now we can only
report relative reductions.
(iv)~At $k\in\{1,2\}$ the redundancy that RR exploits collapses to
zero on Erd\H{o}s--R\'enyi inputs; an LS variant tailored to the
small-radius regime — perhaps with $(2,1)$- or $(3,2)$-swaps as
in~\cite{zhu2024} — is the natural next step.
(v)~The BFS and counter-update phases of both RR and LS are
embarrassingly parallel across $v\in D$; a multi-thread
implementation should remove most of the GT bottleneck observed in
Section~\ref{sec:exp:time}.


% =============================================================
% REFERENCES
% =============================================================
\bibliographystyle{plain}
% \bibliography{references}  % Uncomment when using a .bib file

\begin{thebibliography}{99}

\bibitem{abed2022hybrid}
S.A.~Abed, H.M.~Rais, J.~Watada, and N.R.~Sabar.
\newblock A hybrid local search algorithm for minimum dominating set
  problems.
\newblock \emph{Engineering Applications of Artificial Intelligence},
  114:105053, 2022.
\newblock \url{https://doi.org/10.1016/j.engappai.2022.105053}

\bibitem{abed2025}
S.A.~Abed, M.~Tareq, A.A.~Mahmood, and F.~Hassan.
\newblock An effective binary equilibrium optimizer based on
  fitness-distance balance for the minimum $k$-dominating set
  problem.
\newblock In \emph{2025 IEEE Int.\ Conf.\ on Artificial Intelligence
  in Engineering and Technology (IICAIET)}, pages 425--430, 2025.
\newblock \url{https://doi.org/10.1109/IICAIET67254.2025.11265404}

\bibitem{alipour2021}
S.~Alipour and E.~Salari.
\newblock On distributed algorithms for minimum dominating set
  problem and beyond.
\newblock \emph{arXiv preprint arXiv:2103.08061}, 2021.
\newblock \url{https://arxiv.org/abs/2103.08061}

\bibitem{amis2000}
A.D.~Amis, R.~Prakash, T.H.P.~Vuong, and D.T.~Huynh.
\newblock Max-min {D}-cluster formation in wireless ad hoc networks.
\newblock In \emph{IEEE INFOCOM}, pages 32--41, 2000.

\bibitem{barrena2025}
E.~Barrena, S.~Bermudo, A.G.~Hern{\'a}ndez-D{\'\i}az,
  A.D.~L{\'o}pez-S{\'a}nchez, and J.A.~Zamudio.
\newblock Finding the minimum $k$-weighted dominating sets using
  heuristic algorithms.
\newblock \emph{Mathematics and Computers in Simulation},
  228:485--497, 2025.
\newblock \url{https://doi.org/10.1016/j.matcom.2024.09.010}

\bibitem{boldi2011}
P.~Boldi, M.~Rosa, and S.~Vigna.
\newblock {HyperANF}: Approximating the neighbourhood function of very
  large graphs on a budget.
\newblock In \emph{Proc.\ 20th Int.\ Conf.\ World Wide Web (WWW)},
  pages 625--634, 2011.

\bibitem{campan2016}
A.~Campan, T.M.~Truta, and M.~Beckerich.
\newblock Fast dominating set algorithms for social networks.
\newblock In \emph{Proc.\ Modern Artificial Intelligence and
  Cognitive Science Conf.\ (MAICS)}, pages 55--62, 2016.

\bibitem{cappart2023}
Q.~Cappart, D.~Ch{\'e}telat, E.B.~Khalil, A.~Lodi, C.~Morris, and
  P.~Veli{\v c}kovi{\'c}.
\newblock Combinatorial optimization and reasoning with graph neural
  networks.
\newblock \emph{Journal of Machine Learning Research}, 24(130):1--61,
  2023.
\newblock \url{https://arxiv.org/abs/2102.09544}

\bibitem{chacon2022}
C.~Chac{\'o}n Sartori and C.~Blum.
\newblock Boosting a genetic algorithm with graph neural networks for
  multi-hop influence maximization in social networks.
\newblock In \emph{Proc.\ FedCSIS}, 2022.

\bibitem{chen2023}
H.~Chen, R.~Chitnis, P.~Eades, and A.~Wirth.
\newblock Sublinear-space streaming algorithms for estimating graph
  parameters on sparse graphs.
\newblock In \emph{Proc.\ Workshop on Algorithms and Data Structures
  (WADS)}, 2023.
\newblock \url{https://arxiv.org/abs/2305.16815}

\bibitem{chen2023deepopt}
J.~Chen, S.~Cai, Y.~Wang, W.~Xu, J.~Ji, and M.~Yin.
\newblock Improved local search for the minimum weight dominating set
  problem in massive graphs by using a deep optimization mechanism.
\newblock \emph{Artificial Intelligence}, 314:103819, 2023.
\newblock \url{https://doi.org/10.1016/j.artint.2022.103819}

\bibitem{cohen2015}
E.~Cohen.
\newblock All-distances sketches, revisited: HIP estimators for massive
  graphs analysis.
\newblock \emph{IEEE Transactions on Knowledge and Data Engineering},
  27(9):2320--2334, 2015.
\newblock Earlier version: arXiv:1306.3284.

\bibitem{dijkstra2024}
T.~Dijkstra, A.~Gagarin, P.~Corcoran, and R.~Lewis.
\newblock Greedy and randomized heuristics for optimization of
  $k$-domination models in digraphs and road networks.
\newblock \emph{arXiv preprint arXiv:2409.04226}, 2024.

\bibitem{garey1979}
M.R.~Garey and D.S.~Johnson.
\newblock \emph{Computers and Intractability: A Guide to the Theory of
  NP-Completeness}.
\newblock W.H.~Freeman, 1979.

\bibitem{grandoni2006}
F.~Grandoni.
\newblock A note on the complexity of minimum dominating set.
\newblock \emph{Journal of Discrete Algorithms}, 4(2):209--214, 2006.

\bibitem{hassin1991}
R.~Hassin and A.~Tamir.
\newblock On the minimum diameter spanning tree problem.
\newblock \emph{Information Processing Letters}, 38(2):109--111, 1991.

\bibitem{haynes1998}
T.W.~Haynes, S.T.~Hedetniemi, and P.J.~Slater.
\newblock \emph{Fundamentals of Domination in Graphs}.
\newblock Marcel Dekker, 1998.

\bibitem{kempe2003}
D.~Kempe, J.~Kleinberg, and \'{E}.~Tardos.
\newblock Maximizing the spread of influence through a social network.
\newblock In \emph{KDD '03}, pages 137--146, 2003.

\bibitem{li2020}
Y.~Li, A.~Potru, and F.~Shahrokhi.
\newblock A performance study of some approximation algorithms for
  minimum dominating set in a graph.
\newblock \emph{arXiv preprint arXiv:2009.04636}, 2020.

\bibitem{li2022restart}
R.~Li, S.~Liu, F.~Wang, J.~Gao, H.~Liu, S.~Hu, and M.~Yin.
\newblock A restart local search algorithm with relaxed configuration
  checking strategy for the minimum $k$-dominating set problem.
\newblock \emph{Knowledge-Based Systems}, 254:109619, 2022.
\newblock \url{https://doi.org/10.1016/j.knosys.2022.109619}

\bibitem{nguyen2020}
T.H.~Nguyen et~al.
\newblock Solving the $k$-dominating set problem on very large-scale
  networks.
\newblock \emph{Computational Social Networks}, 7:4, 2020.
\newblock \url{https://doi.org/10.1186/s40649-020-00078-5}

\bibitem{ni2021}
R.~Ni, X.~Li, F.~Li, X.~Gao, and G.~Chen.
\newblock {FastCover}: An unsupervised learning framework for
  multi-hop influence maximization in social networks.
\newblock \emph{arXiv preprint arXiv:2111.00463}, 2021.

\bibitem{parraInza2024EJOR}
E.~Parra Inza, N.~Vakhania, J.M.~Sigarreta, and J.A.~Hern{\'a}ndez Mira.
\newblock Exact and heuristic algorithms for the domination problem.
\newblock \emph{European Journal of Operational Research},
  313(3):926--936, 2024.
\newblock \url{https://arxiv.org/abs/2211.07019}

\bibitem{parraInza2024MDPI}
E.~Parra Inza, J.M.~Sigarreta, and N.~Vakhania.
\newblock Reducing dominating sets in graphs.
\newblock \emph{Algorithms (MDPI)}, 17(6):258, 2024.
\newblock \url{https://arxiv.org/abs/2404.18028}

\bibitem{raz1997}
R.~Raz and S.~Safra.
\newblock A sub-constant error-probability low-degree test, and a
  sub-constant error-probability {PCP} characterization of {NP}.
\newblock In \emph{STOC '97}, pages 475--484, 1997.

\bibitem{rooij2011}
J.M.M.~van Rooij and H.L.~Bodlaender.
\newblock Exact algorithms for dominating set.
\newblock \emph{Discrete Applied Mathematics}, 159(17):2147--2164,
  2011.

\bibitem{xiong2024}
Y.~Xiong and M.~Xiao.
\newblock Exactly solving minimum dominating set and its
  generalization.
\newblock In \emph{Proc.\ Int.\ Joint Conf.\ on Artificial Intelligence
  (IJCAI)}, 2024.
\newblock \url{https://www.ijcai.org/proceedings/2024/780}

\bibitem{zhu2024}
E.~Zhu, Y.~Zhang, S.~Wang, D.~Strash, and C.~Liu.
\newblock A dual-mode local search algorithm for solving the minimum
  dominating set problem.
\newblock \emph{Knowledge-Based Systems}, 298:111950, 2024.
\newblock \url{https://doi.org/10.1016/j.knosys.2024.111950}

\end{thebibliography}


% =============================================================
% APPENDIX
% =============================================================
\appendix

\section{Pseudocode Details}

\subsection{$k$-hop Neighborhood Computation}

\begin{algorithm}[ht]
\caption{$N_k(G,v,k)$ — BFS $k$ hops}
\label{alg:nk}
\begin{algorithmic}[1]
\Require Graph $G$, vertex $v$, integer $k$
\Ensure $N_k(v)=\{u\in V:\mathrm{dist}(u,v)\le k\}$
\State $\text{visited}\gets\{v\}$;\; $\text{frontier}\gets\{v\}$
\For{$i=1$ \textbf{to} $k$}
  \State $\text{next}\gets\emptyset$
  \For{each $u\in\text{frontier}$}
    \State $\text{next}\gets\text{next}\cup(N(u)\setminus\text{visited})$
  \EndFor
  \State $\text{visited}\gets\text{visited}\cup\text{next}$
  \State $\text{frontier}\gets\text{next}$
\EndFor
\State \Return visited
\end{algorithmic}
\end{algorithm}


\section{Per-Graph Timing Breakdown for GT$+$RR}
\label{app:per-graph-time}

Table~\ref{tab:per-graph-rr-time} reports the GT and GT$+$RR running
times per graph at $k=3$. The RR-only overhead is the difference of
consecutive rows.

\begin{table}[ht]
\centering
\caption{Per-graph wall-clock time for GT and GT$+$RR ($k=3$). The RR
  overhead column is $\text{Time(GT$+$RR)}-\text{Time(GT)}$.}
\label{tab:per-graph-rr-time}
\small
\begin{tabular}{lrrrr}
\toprule
\textbf{Graph} & \textbf{GT (s)} & \textbf{GT$+$RR (s)}
  & \textbf{RR overhead (s)} & \textbf{GT$+$RR$+$LS total (s)}\\
\midrule
rand\_100k\_1 & 28.56 & 29.05 & 0.49 & 172.89 \\
rand\_100k\_2 & 29.73 & 30.05 & 0.32 & 169.37 \\
rand\_100k\_3 & 29.20 & 31.72 & 2.52 & 180.14 \\
rand\_100k\_4 & 27.92 & 28.50 & 0.58 & 186.42 \\
rand\_100k\_5 & 32.29 & 36.32 & 4.03 & 191.95 \\
rand\_100k\_6 & 31.69 & 33.40 & 1.71 & 188.85 \\
rand\_100k\_7 & 31.25 & 33.64 & 2.39 & 196.77 \\
rand\_100k\_8 & 32.53 & 32.96 & 0.43 & 193.19 \\
rand\_100k\_9 & 32.05 & 32.58 & 0.53 & 191.02 \\
\midrule
\textbf{Average} & \textbf{30.58} & \textbf{32.02}
  & \textbf{1.44} & \textbf{185.62} \\
\bottomrule
\end{tabular}
\end{table}

The RR overhead averages $1.44$\,s and reaches $4.03$\,s in the
worst case (rand\_100k\_5), consistent with the per-graph
characterization stated in the abstract. The few outliers
(rand\_100k\_3, rand\_100k\_5) likely correspond to runs where Phase~1
or Phase~2 of RR competed for the same memory-cache lines as the
preceding GT phase; the asymptotic $O(|D|\cdot\Nk)$ bound is unchanged.


\section{LocalSwap Activity per Graph}
\label{app:per-graph-ls}

Table~\ref{tab:per-graph-ls} reports the LS-side counters for the
GT$+$LS$+$RR pipeline at $k=3$. ``Removed'' is the number of
\mbox{Case-A} (redundant) deletions; ``Swaps'' is the number of
$(1,1)$-swaps; ``Passes'' is the number of LS passes executed before
the algorithm terminated.

\begin{table}[ht]
\centering
\caption{LocalSwap activity within the GT$+$LS$+$RR pipeline,
  $k=3$. Maximum passes $P=10$.}
\label{tab:per-graph-ls}
\begin{tabular}{lrrr}
\toprule
\textbf{Graph} & \textbf{Removed} & \textbf{Swaps} & \textbf{Passes}\\
\midrule
rand\_100k\_1 & 1{,}098 & 5{,}994 & 10\\
rand\_100k\_2 & 1{,}125 & 5{,}846 & 10\\
rand\_100k\_3 & 1{,}103 & 6{,}262 & 10\\
rand\_100k\_4 & 1{,}120 & 5{,}799 & 10\\
rand\_100k\_5 & 1{,}083 & 5{,}887 & 10\\
rand\_100k\_6 & 1{,}156 & 5{,}855 & 10\\
rand\_100k\_7 & 1{,}068 & 5{,}824 & 10\\
rand\_100k\_8 & 1{,}055 & 5{,}818 & 10\\
rand\_100k\_9 & 1{,}137 & 6{,}133 & 10\\
\midrule
\textbf{Average} & \textbf{1{,}105} & \textbf{5{,}935} & \textbf{10}\\
\bottomrule
\end{tabular}
\end{table}

Every run reaches the pass cap $P=10$, suggesting LS would continue to
make small improvements with more passes. Each pass executes
${\sim}600$ swaps and removes ${\sim}110$ vertices on average, broadly
matching the per-pass cost analysis in
Theorem~\ref{thm:ls-complexity}.


\subsection{Heuristic Sweep on \texttt{rand\_100k\_2}}

\begin{table}[ht]
\centering
\caption{All construction heuristics across $k=1,2,3$ on
  \texttt{rand\_100k\_2}.}
\label{tab:full-k}
\begin{tabular}{crrrrr}
\toprule
$k$ & \textbf{GD} & \textbf{GT} & \textbf{KC} & \textbf{TH}
   & \textbf{GT$+$RR}\\
\midrule
1 & 30{,}741 & 28{,}441 & 30{,}753 & 30{,}741 & 28{,}441 \\
2 & 10{,}634 &  9{,}853 & 10{,}654 & 10{,}692 &  9{,}853 \\
3 &  4{,}921 &  3{,}700 &  4{,}926 &  4{,}932 &  3{,}303 \\
\bottomrule
\end{tabular}
\end{table}


\end{document}
